\documentclass[10pt,conference]{IEEEtran}
% Research article on dsp-props. Every number comes from numbers.tex and the
% tables/*.tex files, generated by make_numbers.py from the result files of
% the repository's scripts (see the "Reproducing" paragraph of Section V).

\usepackage[T1]{fontenc}
\usepackage[utf8]{inputenc}
\usepackage{amsmath,amssymb,amsthm}
\usepackage{booktabs}
\usepackage{graphicx}
\usepackage{xcolor}
\usepackage{url}
\usepackage[hidelinks]{hyperref}
\usepackage{cite}
\usepackage{pgfplots}
\pgfplotsset{compat=1.18}

\newtheorem{proposition}{Proposition}
\newcommand{\eps}{\varepsilon}
\newcommand{\DFT}{\mathcal{F}}
\newcommand{\code}[1]{\texttt{#1}}
\newcommand{\ie}{i.e.,\ }
\newcommand{\eg}{e.g.,\ }

% Generated by make_numbers.py; do not edit.
\newcommand{\ablationDisagree}{0}
\newcommand{\campaignTargets}{46}
\newcommand{\campaignTriaged}{three}
\newcommand{\campaignTriagedShare}{6.5\,\%}
\newcommand{\cexMedDsp}{3}
\newcommand{\cexMedFix}{3}
\newcommand{\cexMedGen}{3}
\newcommand{\cexMedUni}{3}
\newcommand{\cexPnnDsp}{24}
\newcommand{\cexPnnFix}{33}
\newcommand{\cexPnnGen}{23}
\newcommand{\cexPnnUni}{17}
\newcommand{\dspKillAssert}{13}
\newcommand{\dspKillGroup}{7}
\newcommand{\dspKillOther}{80}
\newcommand{\dspKillRelation}{69}
\newcommand{\dspKillRelationStable}{in every seed}
\newcommand{\dspKillTimeout}{4}
\newcommand{\dspKillTolerance}{56}
\newcommand{\dspKillsSeedZero}{149}
\newcommand{\envDescription}{Python 3.13.16, NumPy 2.5.3, SciPy 1.18.1, Hypothesis 6.170.2, on a 4-core Linux x86-64 machine}
\newcommand{\firstMedDsp}{3}
\newcommand{\firstMedFix}{3}
\newcommand{\firstMedGen}{2}
\newcommand{\firstMedUni}{2}
\newcommand{\firstPnnDsp}{8}
\newcommand{\firstPnnFix}{8}
\newcommand{\firstPnnGen}{21}
\newcommand{\firstPnnUni}{21}
\newcommand{\fpComparisons}{248\,000}
\newcommand{\fpFailed}{no}
\newcommand{\fpRuns}{2\,000}
\newcommand{\fpSeeds}{20}
\newcommand{\genCloseBelowOne}{99\,\%}
\newcommand{\genEmptyBlock}{66\,\%}
\newcommand{\genOnGrid}{41\,\%}
\newcommand{\genOneBlock}{44\,\%}
\newcommand{\genPathoCount}{19}
\newcommand{\genPathoLong}{7\,\%}
\newcommand{\genPathoLongDraws}{531}
\newcommand{\genPathoRuns}{20}
\newcommand{\genPathoShare}{15\,\%}
\newcommand{\genPathoShareMax}{16\,\%}
\newcommand{\genPathoShareMin}{14\,\%}
\newcommand{\genZeroNyq}{23\,\%}
\newcommand{\handDspAlways}{31}
\newcommand{\handDspCexPnn}{19}
\newcommand{\handDspFalseAlarms}{0}
\newcommand{\handDspFirstMed}{3}
\newcommand{\handDspFirstPnn}{8}
\newcommand{\handDspNever}{0}
\newcommand{\handFirstThree}{26}
\newcommand{\handFixAlways}{30}
\newcommand{\handFixCexPnn}{57}
\newcommand{\handFixFalseAlarms}{0}
\newcommand{\handFixFirstMed}{3}
\newcommand{\handFixFirstPnn}{13}
\newcommand{\handFixNever}{0}
\newcommand{\handGenAlways}{31}
\newcommand{\handGenCexPnn}{36}
\newcommand{\handGenFalseAlarms}{0}
\newcommand{\handGenFirstMed}{8}
\newcommand{\handGenFirstPnn}{27}
\newcommand{\handGenNever}{0}
\newcommand{\handKilled}{all}
\newcommand{\handMedianLen}{2}
\newcommand{\handUniAlways}{31}
\newcommand{\handUniCexPnn}{34}
\newcommand{\handUniFalseAlarms}{0}
\newcommand{\handUniFirstMed}{8}
\newcommand{\handUniFirstPnn}{20}
\newcommand{\handUniNever}{0}
\newcommand{\killAccDsp}{--}
\newcommand{\killAccFix}{--}
\newcommand{\killAccGen}{--}
\newcommand{\killAccSop}{--}
\newcommand{\killAccUni}{--}
\newcommand{\killDspOnly}{6}
\newcommand{\killFuncDsp}{90\,\%}
\newcommand{\killFuncFix}{90\,\%}
\newcommand{\killFuncGen}{89\,\%}
\newcommand{\killFuncSop}{96\,\%}
\newcommand{\killFuncUni}{89\,\%}
\newcommand{\killNeither}{0}
\newcommand{\killQuietDsp}{87\,\%}
\newcommand{\killQuietFix}{87\,\%}
\newcommand{\killQuietGen}{86\,\% (84--87)}
\newcommand{\killQuietSop}{95\,\%}
\newcommand{\killQuietUni}{86\,\% (84--87)}
\newcommand{\killSopOnly}{17}
\newcommand{\killUnion}{100\,\%}
\newcommand{\knownBugFlagged}{5}
\newcommand{\knownBugRuns}{5}
\newcommand{\marginExact}{63\,\%}
\newcommand{\marginFactor}{5}
\newcommand{\marginMax}{$0.21$}
\newcommand{\marginPnn}{$0.038$}
\newcommand{\nAccuracy}{0}
\newcommand{\nEquivalent}{17}
\newcommand{\nFuncLoud}{83}
\newcommand{\nFuncQuiet}{83}
\newcommand{\nFunctional}{166}
\newcommand{\nMutAuto}{183}
\newcommand{\nMutFft}{76}
\newcommand{\nMutHand}{31}
\newcommand{\nMutLfilter}{55}
\newcommand{\nMutOla}{16}
\newcommand{\nMutUpfirdn}{36}
\newcommand{\nReclassified}{8}
\newcommand{\nSeeds}{5}
\newcommand{\numpyMaxEps}{4.1}
\newcommand{\oracleFirstEquivalent}{25}
\newcommand{\oracleFirstFunctional}{158}
\newcommand{\reclDspAllVariants}{every}
\newcommand{\reclDspOnly}{4}
\newcommand{\reclDspOnlySubject}{upfirdn}
\newcommand{\reclSopOnly}{4}
\newcommand{\reclSopOnlySubject}{lfilter}
\newcommand{\refDftMax}{$1.0\times10^{-3}$}
\newcommand{\refLfilterMax}{$5.4\times10^{-3}$}
\newcommand{\scipyVersion}{1.18.1}
\newcommand{\secMedDsp}{5.61}
\newcommand{\secMedFix}{5.56}
\newcommand{\secMedGen}{3.42}
\newcommand{\secMedSop}{0.001}
\newcommand{\secMedUni}{3.4}
\newcommand{\surveyHardcodedAll}{220}
\newcommand{\surveyLibrosaCompared}{358}
\newcommand{\surveyLibrosaTests}{916}
\newcommand{\surveyScipyAbsCount}{323}
\newcommand{\surveyScipyAbsDistinct}{37}
\newcommand{\surveyScipyComparisons}{1\,645}
\newcommand{\surveyScipyDistinctTol}{66}
\newcommand{\surveyScipyEmpty}{62}
\newcommand{\surveyScipyExplicit}{584}
\newcommand{\surveyScipyHardcoded}{169}
\newcommand{\surveyScipyRelBelowEps}{9}
\newcommand{\surveyScipyRelCount}{240}
\newcommand{\surveyScipyRelDecades}{14}
\newcommand{\surveyScipyRelDistinct}{26}
\newcommand{\surveyScipyRelMax}{$0.1$}
\newcommand{\surveyScipyRelMin}{$10^{-15}$}
\newcommand{\surveyScipyTests}{1\,095}
\newcommand{\surveyScipyTolDecades}{31}
\newcommand{\surveyScipyTolMax}{$3.01$}
\newcommand{\surveyScipyTolMin}{$10^{-31}$}
\newcommand{\surveyScipyTopTols}{$10^{-13}$ (66), $10^{-7}$ (31), $10^{-15}$ (30)}
\newcommand{\surveyScipyZi}{19}
\newcommand{\surveyTotalTests}{2\,388}
\newcommand{\surveyZiEmpty}{2}
\newcommand{\timeoutSeconds}{90}
\newcommand{\torchMaxEps}{195}
\newcommand{\torchWorstN}{596}

% Analysis of the results. The numbers are macros from numbers.tex (make_numbers.py).
\newcommand{\rqOneText}{dsp-props detects \killFuncDsp{} of the \nFunctional{}
functional mutants and the ported suites \killFuncSop{}; neither raises an
alarm on an equivalent mutant, and of the ablations only generic PBT does
(Section~\ref{sec:rq2}). The two sets of detected faults are not nested:
\killSopOnly{} mutants are detected only by the suites, \killDspOnly{} only by
dsp-props, and together they detect \killUnion{} of them.

Every miss of dsp-props has the same cause: the fault shows only for operator
parameters that no declaration uses. A declaration fixes the coefficients of a
filter or the rates of a resampler, and dsp-props then varies the input signal.
Seven \code{lfilter} mutants break filters of order 0 or 1, or filters whose
numerator is shorter than the denominator, and the declared filters have orders
4 and 2 with numerators at least as long. The \code{upfirdn} survivors change the
length of the full output, which the declarations never see since they reach
\code{upfirdn} through \code{resample\_poly} and through a FIR filter truncated to
the input length, or ignore the last input sample, which every relation sees
identically on both of its sides. The suites miss the opposite kind of fault.
Their only comparisons with expected values on the FFT use length 30, which the
subject computes with Bluestein's algorithm; on powers of two they only check
that \code{ifft} inverts \code{fft}, which two mutants of the radix-2 twiddle
factors still satisfy. They also always pass the rates of \code{upfirdn}
explicitly, so four mutants of its default arguments survive.

dsp-props therefore does not detect more of these faults than suites refined
over many years; it detects other ones, without expected values, and its misses
are a matter of declarations rather than of relations: declaring a first-order
filter would have caught the \code{lfilter} survivors.}

\newcommand{\rqTwoText}{On the generated mutants the four randomised variants
detect almost the same faults (Table~\ref{tab:comparison}). The exceptions are
three FFT mutants that the variants with uniform inputs miss in some seeds,
while dsp-props detects them in every seed; two of them are the twiddle-factor
mutants that the ported suites miss. Generic PBT also raised a false alarm in
one of its \nSeeds{} seeds, on the correct FFT and on every FFT mutant
equivalent to it: Hypothesis drew a subnormal input of norm $10^{-323}$, for
which a purely relative tolerance is zero and a rounding error of one subnormal
fails Parseval's identity. The error model of dsp-props includes an absolute
floor, and neither dsp-props nor its other ablations raised an alarm. The
other mutants change the output by orders of magnitude, and almost any input
reveals them:
the median number of inputs before the first failure is \firstMedDsp{} for
dsp-props and \firstMedUni{} for uniform inputs, and the median counterexample
has \cexMedDsp{} and \cexMedUni{} samples. The generators and tolerances do
not change what is detected here, and the generators cost time, with a median of
\secMedDsp{}~s per mutant against \secMedUni{}~s for uniform inputs and
\secMedSop{}~s for a ported suite.

The hand-written faults, which include the accuracy faults, separate the
variants (Table~\ref{tab:hand}). Every variant detects every fault in every seed,
except the fixed tolerance, which misses one resampler fault in one seed, but
not equally fast. With the domain generators, the median failure appears at
input \handDspFirstMed{} (90th percentile \handDspFirstPnn{}); with uniform inputs
at input \handUniFirstMed{} (\handUniFirstPnn{}), and with generic PBT at input
\handGenFirstMed{} (\handGenFirstPnn{}). Counterexamples are smaller too: the 90th percentile of their length is
\handDspCexPnn{} samples for dsp-props, \handUniCexPnn{} with uniform inputs, and
\handFixCexPnn{} with a fixed tolerance. No variant raised a false alarm on
the correct references of this table. With 200 inputs per relation, these differences do not change the detection
rate. We expect them to matter with a smaller budget or on faults that need a
specific input, but we did not measure this. One hand-written fault, a filter
whose gain is wrong by 1\,\%, is detected only because its transfer function is
declared: it is still linear, time-invariant and causal.}

\newcommand{\conclusionResult}{On \nFunctional{} generated faults it detects
\killFuncDsp{}, against \killFuncSop{} for the ported test suites of NumPy and
SciPy, and the two are complementary: together they detect \killUnion. Its
misses come from operator parameters that no declaration exercised; the misses
of the suites, from code paths that their few expected values do not reach.}

\newcommand{\reportedUpstream}{We reported the empty-block faults to
SciPy\footnote{\url{https://github.com/scipy/scipy/issues/26406}} and
contributed a fix\footnote{\url{https://github.com/scipy/scipy/pull/26415}},
the dropped delay of \code{tf2sos} to
SciPy\footnote{\url{https://github.com/scipy/scipy/issues/26417}}, and the
clamping to torchaudio\footnote{\url{https://github.com/pytorch/audio/issues/4238}}.}


\begin{document}

\title{Testing Digital Signal Processing Code\\without an Oracle: Metamorphic Relations,\\Domain Generators and Error-Model Tolerances}

\author{\IEEEauthorblockN{Cl\'ement Basdevant}
\IEEEauthorblockA{Independent Researcher\\
ORCID 0009-0004-7496-2365}}

\maketitle

\begin{abstract}
Signal processing code is hard to test: the correct output of a filter, a
transform or a resampler on an arbitrary input is unknown, the inputs that break
such code are specific, and floating-point results can only be compared within
a tolerance. A survey of five DSP libraries shows that their tests compare with
typed-in values, never generate their inputs, and use hand-written tolerances
chosen test by test (\surveyScipyRelDistinct{} distinct relative tolerances in
\code{scipy.signal} alone). We propose a method
that derives the test from a statement of what an operator is: metamorphic
relations proved for transforms, linear time-invariant filters, rational
resamplers and streaming processors replace the expected output; inputs are
drawn from signal families known to break DSP code; tolerances take the form
of forward error bounds, with one constant per operator class instead of one
value per test, a derived bound for linear recursions, and an optional
calibration against extended precision. On \nFunctional{} functional faults
injected by mutation into four small reference implementations, it detects
\killFuncDsp{}, fewer than the ported test suites of NumPy and SciPy
(\killFuncSop{}); the two detect partly different faults and together detect
all of them. More than half of its detections (\dspKillOther{} of
\dspKillsSeedZero) come from the mutant raising an exception or hanging rather
than from a violated relation. Ablations show no
net gain in detection from the domain generators or the error-model
tolerances: they mainly make failures appear after fewer inputs, with smaller
counterexamples on hand-written faults. On the correct implementations and
declarations we tested, the method raised no false alarm in \fpRuns{} runs.
Applied to SciPy, it rediscovered a known, already fixed fault (issue 22571)
and found one new edge-case exception, \code{sosfilt} on an empty input (issue
26406, not yet confirmed by the maintainers).
\end{abstract}

\begin{IEEEkeywords}
software testing, metamorphic testing, property-based testing, digital signal processing, floating-point error, test oracle
\end{IEEEkeywords}

% ===========================================================================
\section{Introduction}
% ===========================================================================

Digital signal processing (DSP) code sits at the bottom of a large software
stack: the fast Fourier transform (FFT), linear filters and sample-rate
converters of NumPy~\cite{harris2020numpy}, SciPy~\cite{virtanen2020scipy} and
PyTorch~\cite{paszke2019pytorch} are called by audio, communication,
geophysics and machine-learning code that rarely re-checks them. A fault at
this level propagates silently: a filter that loses its state on an empty
block does not crash the application that streams audio through it, it makes
the output diverge. Testing these operators faces three problems at once.

\paragraph*{P1, no oracle} For most inputs nobody knows the exact output of a
filter or a transform. This is the \emph{oracle problem} of software
testing~\cite{barr2015oracle}, and it is acute for scientific
software~\cite{kanewala2014scientific}. In practice, tests compare against
values computed by another tool and typed into the test file, or against a
second implementation that may share the same assumptions. In the suites we
surveyed (Section~\ref{sec:survey}), \surveyScipyHardcoded{} tests of
\code{scipy.signal} compare with hard-coded arrays.

\paragraph*{P2, inputs that do not exercise the faults} Hand-written tests use
a few lengths and smooth or random signals. Faults in DSP code, however, hide in
specific places: lengths that are prime, odd or one away from a power of two
(which change the FFT algorithm), tones between two DFT bins or close to the
Nyquist frequency, extreme amplitudes, and, for streaming code, empty or
one-sample blocks. A concrete case: \code{scipy.signal.lfilter} returned an
uninitialised final state when given an empty block and an initial state
(SciPy issue 22571\footnote{\url{https://github.com/scipy/scipy/issues/22571}}),
so that any block-processing loop silently diverged after an empty block, and
\code{sosfilt} still raises on such input. Across the \surveyTotalTests{} tests
we surveyed, only \surveyZiEmpty{} combine an initial state with an empty
input, and none generates its inputs.

\paragraph*{P3, no principled tolerance} Even when the expected output is
known, floating-point results are only close to it, and \emph{how close} depends
on the operation, the length and the conditioning of the problem. Tests
therefore use tolerances chosen by hand: \surveyScipyDistinctTol{} distinct
values in \code{scipy.signal}, among them \surveyScipyRelDistinct{} relative
tolerances that range from \surveyScipyRelMin{} to \surveyScipyRelMax{} once the
values below machine precision, which ask for exact equality, are set aside.
A tolerance that is too loose hides a single-precision intermediate; one that is
too tight produces false alarms, which teach maintainers to loosen tolerances
further.

Property-based testing~\cite{claessen2000quickcheck,maciver2019hypothesis}
answers P2 in general by generating inputs, and metamorphic
testing~\cite{chen2018metamorphic,segura2016metamorphic} answers P1 in general
by checking relations between outputs instead of the outputs themselves.
Neither tells a DSP developer which relations hold for which operator, which
signals to generate, nor how to compare floating-point results. And a method
that claims to find faults in numerical libraries must itself be shown not to
raise false alarms, since a single spurious report costs a maintainer's trust.

This article answers the three problems with one method, implemented as a
tool, and evaluates it with a protocol designed to measure both its power and
its soundness. The evaluation is a negative result as much as a positive one:
on generated faults, the method detects fewer faults than the existing test
suites, and its domain-specific components do not improve detection
measurably. Its contributions are:
\begin{itemize}
\item \emph{A method} that replaces the expected output by metamorphic relations
proved for four classes of DSP operators (transforms, linear time-invariant
filters, rational resamplers, streaming processors), draws inputs from the
signal families that break DSP code, and gives each tolerance the form of a
forward error bound, with one constant per operator class instead of one value
per test (Section~\ref{sec:method}).
\item \emph{An evaluation protocol} that compares the method with the test
suites of NumPy and SciPy on automatically generated faults, isolates the
contribution of the generators and of the tolerances by ablation, and validates
the method itself: false alarms on correct code, margin between observed error
and tolerance, accuracy of the references against 50-digit arithmetic, and an
audit of the generators (Section~\ref{sec:protocol}).
\item \emph{A campaign on widely used libraries} that rediscovered a known
fault of SciPy and found an edge-case exception, both reported to the
maintainers (Section~\ref{sec:rq4}).
\item \emph{A survey} of how five DSP libraries test themselves, which grounds
P1--P3 in counts rather than impressions (Section~\ref{sec:survey}).
\end{itemize}
The method is implemented in dsp-props, an open-source layer on top of
Hypothesis; the library, the experiments and their data are available at
\url{https://github.com/clement-basdevant/dsp-props} under a BSD licence.

Section~\ref{sec:related} reviews related work, Section~\ref{sec:survey}
reports the survey, Section~\ref{sec:method} describes the method,
Section~\ref{sec:protocol} the evaluation protocol and Section~\ref{sec:results}
its results. Section~\ref{sec:discussion} discusses what the results mean for
practitioners, and Section~\ref{sec:conclusion} concludes.

% ===========================================================================
\section{Background and related work}
\label{sec:related}
% ===========================================================================

\subsection{The oracle problem and metamorphic testing}
A test oracle decides whether an output is correct. When no such procedure
exists, or when it is too costly, testing faces the oracle
problem~\cite{barr2015oracle}; scientific software is the textbook case, since
it is often written precisely because the answer is not known in
advance~\cite{kanewala2014scientific}. Metamorphic
testing~\cite{chen2018metamorphic,segura2016metamorphic} avoids the oracle by
checking \emph{metamorphic relations}: if $f$ is correct, then for related
inputs $x$ and $x'$ the outputs $f(x)$ and $f(x')$ must satisfy a known
relation, \eg $f(x+x') = f(x)+f(x')$ for a linear $f$. It has been applied
among others to bioinformatics~\cite{chen2009bioinformatics} and image
processing~\cite{jafari2024image}. The quality of the approach depends on the
relations: a relation that holds for the faulty program too detects nothing,
and finding relations is the main cost of the method. For DSP operators the
relations are classical theorems (linearity, shift and convolution theorems,
Parseval's identity), but which ones apply depends on what the operator is, and
on finite sequences their exact form at the edges matters.

\subsection{Property-based testing}
Property-based testing (PBT), introduced by
QuickCheck~\cite{claessen2000quickcheck} and available in Python through
Hypothesis~\cite{maciver2019hypothesis}, states a property over all inputs of a
type, generates inputs and, when one fails, \emph{shrinks} it to a smaller,
simpler failing input~\cite{maciver2020reduction}. PBT supplies the machinery to check
metamorphic relations, but not the relations themselves, nor the generators of
meaningful inputs. Interviews of PBT users identify writing properties and
generators as their main difficulty~\cite{goldstein2024pbt}, and Maaz
et~al.\ use language models to propose properties for Python
libraries~\cite{maaz2025agentic}. Generic generators of floating-point arrays
also produce values, such as subnormal numbers, for which a purely relative
tolerance is meaningless; Section~\ref{sec:rq2} shows a false alarm of this
kind.

\subsection{Self-testing of numerical programs}
Blum, Luby and Rubinfeld~\cite{blum1993selftesting} showed that a program
computing a linear function can be tested, and even corrected, by checking
random linear identities. Erg\"un~\cite{ergun1995testing,ergun2000selftesting}
extended the idea to multivariate linear functions such as the DFT, and
Erg\"un, Kumar and Rubinfeld~\cite{ergun2001checking} to computations that are
only approximately correct. The FFTW library~\cite{frigo2005fftw} ships a
self-test that implements Erg\"un's tester and checks linearity, impulse
responses and time and frequency shifts~\cite{fftwverify}. These works establish that
linearity and shift identities characterise the DFT; they do not cover
filters, resamplers or streaming code, and they treat the floating-point error
with a single threshold.

\subsection{Differential testing and numerical faults}
Differential testing compares several implementations of the same
function~\cite{vanover2020fpdiff}. It needs at least two of them, and they may
share an algorithm or an assumption, \eg the same handling of an empty block.
Di Franco et~al.\ studied 269 numerical bugs reported in NumPy, SciPy,
LAPACK, GSL and Elemental~\cite{difranco2017floatingpoint}. They selected bugs
confirmed by developers, most of them closed with a fix, and analysed 183 of
those fixes: the study shows that numerical faults reach widely used, mature
libraries, not that these particular faults remain in them.

\subsection{Testing numerical and deep-learning libraries}
Several tools test deep-learning libraries, whose operators overlap with DSP
code. CRADLE compares the outputs of several backends of the same
model~\cite{pham2019cradle} and Audee generates models and inputs for such
comparisons~\cite{guo2020audee}; both are differential testing. FreeFuzz
mines API calls from open-source code~\cite{wei2022freefuzz}, DeepREL infers
pairs of APIs that should give related results~\cite{deng2022deeprel}, DocTer
extracts input constraints from the documentation~\cite{xie2022docter}, and
TitanFuzz generates programs with a language model~\cite{deng2023titanfuzz}.
These tools find crashes and inconsistencies across a whole API with little
domain knowledge, but do not aim at deriving the tolerance of numerical
comparisons from the operator; our
method is narrower and states, for a few operator classes, which relations
hold and how large their error may be. Kanewala and Bieman survey the testing
of scientific software and name the oracle problem and the tolerance on
floating-point results among its main
difficulties~\cite{kanewala2014scientific}.

\subsection{Tolerances of test assertions}
Dutta et~al.\ found that many flaky tests in probabilistic and
machine-learning projects come from assertions whose bounds do not account for
randomness, and proposed FLASH, which detects them by running the tests with
many seeds~\cite{dutta2020flaky}. FLEX repairs such assertions by fitting an
extreme-value distribution to the values observed over many runs and setting
the bound to a chosen percentile~\cite{dutta2021flex}, and TERA reduces the cost of stochastic
regression tests by tuning their parameters~\cite{dutta2021tera}.
FLEX and our calibration share a principle, measuring the variability of a
correct program to set a bound, but FLEX targets randomness from seeds while we
target rounding error, and our default tolerances come from the form of an
error bound rather than from observations.

\subsection{Floating-point error analysis}
Static analysers bound or reduce the rounding error of a given expression or
program: Herbie rewrites expressions to improve their
accuracy~\cite{panchekha2015herbie}, FPTaylor bounds round-off errors with
symbolic Taylor expansions~\cite{solovyev2019fptaylor}, and Rosa and Daisy
compile and analyse programs over the reals with sound error
bounds~\cite{darulova2014rosa,darulova2018daisy}. Dynamic tools search for
inputs that cause large errors, guided by the condition of each
operation~\cite{zou2020atomu} or by symbolic
execution~\cite{guo2020fpgen}. These tools analyse the implementation, which
our method treats as a black box; their bounds could replace our
engineering constants for code they can analyse. Linear filters have their
own line of work: abstract interpretation bounds the values of digital
filters~\cite{feret2004filters,monniaux2005filters}, and work on fixed-point
implementations evaluates the worst-case peak gain of a filter with
guaranteed accuracy and uses it to choose number
formats~\cite{volkova2015wcpg,lopez2014formatting}. The bound of
Appendix~\ref{app:iir} uses the same quantity.

\subsection{Floating-point error of DSP operators}
The FFT computed by a Cooley--Tukey algorithm has a worst-case relative error
bounded by $O(\eps\log N)$~\cite{higham2002accuracy}, where
$\eps=2^{-52}$ is the machine epsilon of double precision
(\code{np.finfo(float).eps}), twice the unit roundoff $u$ in which such bounds
are usually stated; we use $\eps$ throughout because the library does. Its
root-mean-square error on random inputs grows only as
$O(\eps\sqrt{\log N})$~\cite{schatzman1996accuracy}; a direct DFT, computed as $N$
dot products of length $N$, has a bound that grows as
$O(\eps N)$~\cite{higham2002accuracy}. For a linear recursion, rounding errors
are fed back through the poles: Appendix~\ref{app:iir} derives from the
standard model of floating-point arithmetic~\cite{higham2002accuracy} a bound
that grows with the $\ell_1$ norm of the impulse response, which is large when
poles approach the unit circle. These bounds are
rarely used in tests, which fall back on fixed relative tolerances; we use them
to set tolerances (Section~\ref{sec:tolerances}).

\subsection{Mutation analysis}
Mutation analysis~\cite{jia2011mutation,papadakis2019mutation} measures the
fault-detection power of a test approach by injecting small syntactic changes,
\emph{mutants}, into a correct program and counting how many the tests detect
(\emph{kill}). Mutation operators replace an arithmetic, relational or bitwise
operator, change a constant, delete a statement, and so on. Offutt
et~al.~\cite{offutt1996mutation} showed that a small selective set of operators
(arithmetic, relational and logical operator replacement, operand replacement
and unary operator insertion) yields almost the same mutation score as the full
set, which keeps mutation affordable. Two
assumptions underpin the method: the \emph{competent programmer} hypothesis,
under which real faults are small deviations from a correct program, and the
\emph{coupling effect}, under which tests that kill simple faults also kill
complex ones~\cite{jia2011mutation}. Just et~al.\ found that mutant detection
correlates with the detection of real faults in Java
programs~\cite{just2014mutants}.

Two difficulties matter here. First, some mutants are \emph{equivalent}: they
compute the same function as the original, and no test can kill them; deciding
equivalence is undecidable in general~\cite{jia2011mutation}, so studies use
heuristics. Numerical code adds a third category: a mutant may change the
output only by rounding, \eg by reordering an addition. Such a mutant is
neither equivalent nor functionally wrong, and whether it should be killed
depends on the accuracy the program promises. We classify mutants with an
independent oracle into equivalent, accuracy and functional faults
(Section~\ref{sec:groundtruth}) and report kill rates per class, so that a
false alarm on an equivalent mutant counts against an approach instead of for
it. Second, first-order mutants are mostly easy to kill; we therefore also use
hand-written faults modelled on real DSP bugs, including accuracy faults.

\subsection{Summary}
Table~\ref{tab:positioning} summarises what each family of approaches leaves to
the developer. Every approach answers part of P1--P3; none supplies relations,
generators and tolerances together for the operators of a DSP library, which is
the gap this article addresses.

\begin{table}[t]
\caption{What each approach supplies to the developer: P1, no need for expected values; P2, generated inputs; P3, a principled tolerance. FLEX sets the bound of an existing assertion from the values observed over many runs, which addresses P3 for randomness rather than rounding}
\label{tab:positioning}
\centering\scriptsize
\begin{tabular}{@{}llll@{}}
\toprule
Approach & P1 & P2 & P3 \\
\midrule
Expected-value tests & -- & -- & -- \\
Differential~\cite{vanover2020fpdiff} & needs 2nd impl. & partly & -- \\
DFT self-testing~\cite{ergun1995testing,fftwverify} & DFT only & yes & -- \\
PBT, MT~\cite{maciver2019hypothesis,chen2018metamorphic} & user's relations & generic & -- \\
DL fuzzing~\cite{wei2022freefuzz,deng2022deeprel} & API pairs & API-wide & -- \\
FLEX~\cite{dutta2021flex} & -- & -- & observed \\
This work & yes & DSP families & yes \\
\bottomrule
\end{tabular}
\end{table}

% ===========================================================================
\section{How DSP libraries test themselves}
\label{sec:survey}
% ===========================================================================

To ground P1--P3 in data, we analysed the test suites of \code{numpy.fft},
\code{scipy.fft}, \code{scipy.signal}, librosa~\cite{mcfee2015librosa} and the
functional module of torchaudio, \surveyTotalTests{} test functions in total.

\paragraph*{Method} A script parses every test file into a syntax tree and
counts, per library: the test functions and those decorated with Hypothesis'
\code{@given}; the numerical comparisons, \ie calls to \code{assert\_allclose},
\code{assert\_array\_almost\_equal}, \code{allclose},
\code{torch.testing.assert\_close} and their relatives; among them, those with a
hand-written tolerance, \ie a literal passed as \code{rtol}, \code{atol},
\code{decimal}, \code{nulp} or similar, converted to a magnitude
($\code{decimal}=d$ means $1.5\times10^{-d}$); the tests that pass an empty
signal, an initial filter state (\code{zi=}), or both; and the tests that compare
with a typed-in array of at least four floating-point literals. These are
syntactic heuristics: a tolerance computed at run time or an expected array
loaded from a file is not counted, so the counts are lower bounds.

\begin{table*}[t]
\caption{Test suites of five DSP libraries (syntactic counts). Tol.: comparisons with a hand-written tolerance; rel.: hand-written relative tolerances (\code{rtol}, \code{nulp} and similar), their number of distinct values and their range (absolute tolerances, in the units of the data, are not comparable across tests and are left out); empty: tests with an empty input; zi+empty: tests that also pass an initial state; hard-coded: tests that compare with an array of at least four typed-in floating-point values}
\label{tab:survey}
\centering\footnotesize
\begin{tabular}{@{}lrrrrrrlrrr@{}}
\toprule
Suite & Tests & Hypothesis & Comparisons & Tol. & Rel. & Distinct & Rel.\ range & Empty & zi+empty & Hard-coded \\
\midrule
NumPy (numpy.fft) & 52 & 0 & 102 & 66 & 2 & 1 & $10^{-14}$ & 1 & 0 & 1 \\
SciPy (scipy.fft) & 96 & 0 & 132 & 11 & 9 & 4 & $10^{-7}$--$3.0\times10^{-6}$ & 0 & 0 & 4 \\
SciPy (scipy.signal) & 1095 & 0 & 1645 & 584 & 240 & 26 & $10^{-31}$--$0.1$ & 62 & 2 & 169 \\
librosa & 916 & 0 & 664 & 60 & 19 & 6 & $10^{-12}$--$0.01$ & 4 & 0 & 36 \\
torchaudio (functional) & 229 & 0 & 13 & 12 & 12 & 6 & $10^{-7}$--$10^{-3}$ & 0 & 0 & 10 \\
\bottomrule
\end{tabular}
\end{table*}


\paragraph*{Three typical tests} The ported suites of Section~\ref{sec:protocol}
give a concrete picture. SciPy's \code{test\_rank\_1\_IIR} filters the ramp
$0,1,\dots,5$ with $b=[1,-1]$, $a=[0.5,-0.5]$ and compares the output with the six
typed-in values $0,2,4,6,8,10$ to six decimals: one input, one filter, one
tolerance. \code{test\_length\_factors} for \code{upfirdn} checks five
combinations of lengths and rates against output vectors whose comment says they
come from MATLAB. \code{test\_vs\_convolve} compares \code{upfirdn} with a
straightforward implementation on 10\,000 random samples, for three
decimation factors, at a relative and absolute tolerance of $10^{-7}$, which is a
form of differential testing within the library. Several tests are regression
tests added after an issue was reported, such as the two \code{test\_short\_x} tests for a
one-sample input with an initial state: edge cases enter the suites after a user
has met them.


\paragraph*{No oracle (P1)} Table~\ref{tab:survey} shows that hard-coded
expected arrays are common: \surveyHardcodedAll{} tests in total, of which
\surveyScipyHardcoded{} in \code{scipy.signal}. Such arrays come from another
tool, MATLAB according to the comments of some tests, or from an earlier run of
the library, and they test the library against that source on a handful of
inputs. The other tests compare
two computations of the same quantity within the library, \eg \code{lfilter}
against \code{convolve}, which is a form of differential testing.

\paragraph*{Few, fixed inputs (P2)} No test in the five suites uses Hypothesis
or any other input generator. In the tests we ported, inputs are written by
hand: short ramps, a few random arrays with a fixed seed, constant signals. Edge cases are present
but scattered: \code{scipy.signal} has \surveyScipyEmpty{} tests with an empty
input and \surveyScipyZi{} with an initial state, but only \surveyZiEmpty{} that
combine both, which is the case behind the faults of Section~\ref{sec:rq4}.

\paragraph*{Hand-written tolerances (P3)} \code{scipy.signal} alone performs
\surveyScipyComparisons{} comparisons, \surveyScipyExplicit{} of them with a
hand-written tolerance, which take \surveyScipyDistinctTol{} distinct values.
Relative and absolute tolerances must be read separately: an absolute
tolerance is in the units of the data (some are in decibels), so its magnitude
says nothing by itself. The \surveyScipyAbsCount{} absolute tolerances take
\surveyScipyAbsDistinct{} values that we do not compare. The
\surveyScipyRelCount{} relative tolerances take \surveyScipyRelDistinct{}
values (Fig.~\ref{fig:tolhist}); \surveyScipyRelBelowEps{} of them lie below
machine precision and ask for exact equality, and the others range from
\surveyScipyRelMin{} to \surveyScipyRelMax, \surveyScipyRelDecades{} orders of
magnitude. The most frequent are \surveyScipyTopTols. A relative tolerance near
$10^{-15}$ asserts agreement to the last bits, and one near $10^{-4}$ would
accept a single-precision computation. The survey does not tell why each value
was chosen; it shows that the tolerance is set test by test rather than
derived from the operation, which is P3.

\begin{figure}[t]
\centering
\begin{tikzpicture}
\begin{axis}[ybar, width=\columnwidth, height=0.42\columnwidth, bar width=3.5pt,
  xlabel={Hand-written relative tolerance (decade, $\log_{10}$)}, ylabel={Comparisons}, ylabel near ticks,
  xmin=-32, xmax=1, ymin=0, tick label style={font=\scriptsize},
  label style={font=\footnotesize}, xtick={-30,-25,-20,-15,-10,-5,0}]
\addplot[fill=blue!60!black, draw=none]
  table[x=decade, y=count, col sep=comma]{tables/tolhist.csv};
\end{axis}
\end{tikzpicture}
\caption{Hand-written relative tolerances in the tests of \code{scipy.signal},
grouped by decade. Absolute tolerances are left out.}
\label{fig:tolhist}
\end{figure}

% ===========================================================================
\section{Method}
\label{sec:method}
% ===========================================================================

The method rests on one idea: the user states \emph{what an operator is}, and
everything else follows from that statement. A DFT, a causal LTI filter with a
given transfer function, an $L/M$ resampler or a streaming processor each
satisfy known theorems (the relations), are known to fail on known inputs (the
generators), and have a known forward error (the tolerance). The method is
implemented in dsp-props as a layer on top of Hypothesis, but its ingredients
do not depend on that implementation.

\subsection{Declarations}
A declaration is a decorator on the operator under test. For instance,
\begin{quote}\footnotesize\begin{verbatim}
@dp.lti_filter(transfer=(b, a))
def my_filter(x):
    return scipy.signal.lfilter(b, a, x)

dp.verify(my_filter)
\end{verbatim}\end{quote}
declares a causal linear time-invariant (LTI) filter with transfer function
$B(z)/A(z)$. \code{dp.verify} runs every relation that applies to the declared
class on generated inputs (200 per relation by default) and raises an error
that carries the smallest failing input found by shrinking, the relation that
failed, the observed error and the tolerance. The declarations available are
transforms (complex DFT, DFT of real input, generic linear transforms such as
the DCT), LTI filters (with an optional transfer function and gain bound),
rational resamplers, and streaming or stateful processors; adapters wrap NumPy,
SciPy, PyTorch and C functions called through \code{ctypes}.

A declaration fixes the parameters of the operator: the coefficients of a
filter, the rates of a resampler. The relations then vary the input signal
only. This is a deliberate choice, since parameters such as filter
coefficients must satisfy constraints (stability, a declared transfer function)
that a generic generator cannot know; Section~\ref{sec:rq1} shows its cost.

\subsection{Relations (answer to P1)}
\label{sec:relations}
Table~\ref{tab:relations} lists the relations for each operator class. Each is
a consequence of the declaration~\cite{oppenheim2010dtsp,vaidyanathan1993multirate} and is proved in Appendix~\ref{app:proofs}; we
state them on finite sequences of length $N$, as implementations compute them,
which matters for shifts and edges. For a transform, the shift theorem is the
\emph{circular} one, and the shift is reduced modulo $N$ in integers so that the
twiddle factors of the reference are themselves accurate. For a filter, a delay
is applied by prepending zeros, and the delayed output is compared with the
original output, also delayed, from zero state. For a resampler, a delay of $M$
input samples must delay the output by exactly $L$ samples, and the output
length must be $\lceil NL/M\rceil$. For a streaming processor, the input is cut
into blocks at generated positions, including empty and one-sample blocks, and
the concatenated outputs must equal the one-shot output.

The relations replace the oracle: no expected output is needed, except for the
optional transfer function, which the user states once instead of typing
expected arrays. Without it, an LTI filter is only checked to be linear,
time-invariant and causal, so a filter with the wrong coefficients but the
right structure passes; this blind spot is measured in Section~\ref{sec:rq1}.

\begin{table*}[t]
\caption{Relations checked for each operator class. $T$: transform; $H$: filter; $R$: resampler; $D_s$: delay by $s$ samples with zero filling; $S_s$: circular shift; $\omega=e^{-2\pi i/N}$}
\label{tab:relations}
\centering\footnotesize
\renewcommand{\arraystretch}{1.15}
\begin{tabular}{@{}lll@{}}
\toprule
Class & Relation & Checked property \\
\midrule
Transform & Linearity & $T(\alpha x+\beta y)=\alpha Tx+\beta Ty$ \\
 & Circular shift theorem & $T(S_s x)[k]=\omega^{ks}\,Tx[k]$ \\
 & Parseval's identity & $\|Tx\|_2^2=c\,\|x\|_2^2$, with $c$ declared \\
 & Hermitian symmetry & $x$ real $\Rightarrow Tx[N-k]=\overline{Tx[k]}$; DC and Nyquist bins real \\
 & Inversion & $T^{-1}(Tx)=x$ \\
 & Convolution theorem & $T(x\circledast y)=Tx\cdot Ty$ \\
\midrule
LTI filter & Linearity & $H(\alpha x+\beta y)=\alpha Hx+\beta Hy$, from zero state \\
 & Time invariance & $H(D_s x)=D_s(Hx)$ \\
 & Causality & changing $x[n]$ for $n\ge k$ leaves $y[n]$ unchanged for $n<k$ \\
 & Frequency response & DFT of the impulse response equals $B(e^{i\theta})/A(e^{i\theta})$ on the DFT grid \\
 & Bounded gain & $|x[n]|\le1$ for all $n$ $\Rightarrow$ $|y[n]|\le G$, with $G$ declared \\
\midrule
Resampler $L/M$ & Linearity & as for filters \\
 & Output length & $|Rx|=\lceil NL/M\rceil$ \\
 & Periodic time invariance & $R(D_{Ms}x)=D_{Ls}(Rx)$ \\
 & DC gain & a constant input gives the declared constant away from the edges \\
 & Alias rejection & a tone above the output Nyquist frequency is attenuated, not folded back \\
\midrule
Streaming & Block invariance & generated splits (Sec.~\ref{sec:generators}), with empty and one-sample blocks, give the one-shot output \\
\bottomrule
\end{tabular}
\end{table*}

\subsection{Generators (answer to P2)}
\label{sec:generators}
Signals are drawn from seven families, each aimed at a known class of fault:
impulses (indexing and edge errors), constants (DC gain, normalisation),
sinusoids at arbitrary frequency, including off-grid ones (leakage and twiddle
errors), sinusoids near the Nyquist frequency (aliasing, sign alternation),
pairs of tones less than one DFT bin apart (resolution, accumulated rounding),
sparse signals and white noise. Amplitudes are drawn from
$\{10^{-100},10^{-3},1,10^{3},10^{100}\}$ to reach underflow and overflow paths.
Lengths are biased towards the values that change the code path of an FFT or
a filter: 1, 2, small primes, and $2^k-1$, $2^k$, $2^k+1$, mixed with uniform
lengths up to 512. The bias is moderate: in runs of 200 inputs, the default,
\genPathoShare{} of the drawn lengths are among these values on average
(\genPathoShareMin{} to \genPathoShareMax{} over \genPathoRuns{} seeds), the
rest coming from the uniform branch. Block splits for streaming code mix three kinds: arbitrary
cut points, which give up to $\max(N,16)$ blocks including empty and one-sample
ones, every block of a single sample, and blocks of a fixed drawn size, so that
a processor that loses its state only after many calls is reached.

The families are ordered so that Hypothesis' shrinking converges to short
impulses of unit amplitude, which yields counterexamples a DSP engineer can read
at a glance: on the hand-written faults of Section~\ref{sec:rq2}, the median
counterexample has \handMedianLen{} samples.

\subsection{Tolerances (answer to P3)}
\label{sec:tolerances}
All comparisons use the $\ell_2$ norm of the error, scaled by the norm of a
reference quantity, and a model of the admissible error. Using norms rather
than element-wise comparisons follows the form of the error bounds, which are
normwise, and avoids spurious failures on output samples close to zero.

\emph{FFT.} For transforms,
\begin{equation}
\|\hat y-y\|_2 \le c\,\eps\,\max(\log_2N,1)\,\|y\|_2 + N\,\mathrm{tiny},
\label{eq:fft}
\end{equation}
the form of the forward error bound of the radix-2
FFT~\cite[Thm.~24.2]{higham2002accuracy}, with a constant $c=20$ that covers
the few chained transforms in a relation, and a floor at $N$ times the smallest
normal number $\mathrm{tiny}$. The floor keeps exact zeros strict while
tolerating underflow on signals of amplitude $10^{-100}$ or subnormal values,
where a purely relative tolerance would be zero. The constant is an
engineering choice, not a derived value: it was set to avoid false alarms on
correct libraries, and Section~\ref{sec:rq3} measures the resulting margin.
Filters without a declared transfer function, resamplers and streaming
processors use the same form with $c=1000$, a looser default meant to cover
implementations by direct convolution or overlap-add, whose error grows faster
than $\log N$. This constant is an engineering choice too, not derived from an
error analysis of these operators; the user can replace it by a calibrated
model. In every case
$\eps$ is the machine epsilon of the input (or of a declared precision), not of
the output: an operator that returns single precision for a double-precision
input is reported instead of being granted a single-precision tolerance.

\emph{Linear recursions.} For a filter $y=(B/A)x$ computed by the direct form~I
recursion, Appendix~\ref{app:iir} derives
\begin{equation}
\|\hat y-y\|_2 \le c\,\eps\,k\,\|h_A\|_1\bigl(\|b\|_1+\|a_{1:}\|_1\|h_{B/A}\|_1\bigr)\|x\|_2,
\label{eq:iir}
\end{equation}
to first order in $\eps$, where $k$ is the number of coefficients other than $a_0=1$, $h_A$ and
$h_{B/A}$ are the impulse responses of $1/A$ and $B/A$, and $a_{1:}$ are the
feedback coefficients. The bound grows as poles approach the unit circle, which
is where recursions lose accuracy, and reduces to the error of a dot product for
a FIR filter. It is the only derived tolerance of the library, and it is
derived for the direct form~I recursion only, to first order. Other
structures, such as the transposed direct form~II of \code{scipy.signal.lfilter},
are checked against it empirically (Section~\ref{sec:rq3}), not by proof. The
factor $\|h_A\|_1$ is the worst-case peak gain of $1/A$, a quantity that the
literature on filter implementation evaluates
rigorously~\cite{volkova2015wcpg}.

\emph{Calibration.} An analytic bound cannot describe every library: an FFT
may switch algorithm with the factorisation of $N$, and its error then varies by
two orders of magnitude between neighbouring lengths (Fig.~\ref{fig:profile}).
\code{dp.calibrate} measures the error of an implementation against an
extended-precision reference (a DFT in \code{long double} where that type is
wider than double, a double-double DFT elsewhere, \eg on Windows and on
macOS arm64) for every length, over several random inputs, and
turns the worst error times a margin of 8 into a per-length tolerance. It
refuses to calibrate when the error exceeds what rounding can explain
($10^6\eps$), so that a fault cannot be absorbed into the tolerance.

\subsection{Differences with the state of the art}
Compared with generic PBT, the method supplies the three pieces of domain
knowledge that PBT leaves to the user; Section~\ref{sec:rq2} measures the
contribution of two of them. Compared with self-testing of the
DFT~\cite{ergun1995testing,fftwverify}, it covers filters, resamplers and
streaming code, and its tolerance follows the operator. Compared with
differential testing~\cite{vanover2020fpdiff}, it needs a single implementation.
Compared with the test suites of Section~\ref{sec:survey}, it needs no expected
values and no tolerance chosen test by test.

% ===========================================================================
\section{Evaluation protocol}
\label{sec:protocol}
% ===========================================================================

We answer four research questions (RQ).
\textbf{RQ1}: does the method detect more faults than the state of practice?
\textbf{RQ2}: what do the domain generators and the error-model tolerances contribute?
\textbf{RQ3}: is the method itself correct, \ie free of false alarms, with sound references and generators?
\textbf{RQ4}: does it find faults in real libraries?
RQ1 and RQ2 use injected faults, RQ3 correct code, RQ4 the libraries
themselves.

\subsection{Subjects}
We wrote four small reference implementations with the interface of a library
function: a textbook radix-2 FFT with Bluestein's algorithm for other lengths;
a direct-form~II transposed \code{lfilter} with initial state; an
\code{upfirdn} that computes each output sample directly, as a dot product of
the filter with the non-zero samples of the upsampled input, which a port of
\code{resample\_poly} calls; and a block FIR filter by overlap-add. They are
short enough for exhaustive mutation. Only the \code{lfilter} subject has the
structure of its library counterpart: NumPy's FFT is a mixed-radix
implementation, and SciPy's \code{upfirdn} is a polyphase filter bank in
compiled code. The subjects therefore share the interfaces of the libraries,
and so their test suites, but not their code. All four agree with NumPy or
SciPy to a relative error below $1.6\times10^{-15}$ and pass every suite
below.

\subsection{Faults}
A mutation tool generates every first-order mutant of each subject with
standard operators, following~\cite{offutt1996mutation,jia2011mutation}: arithmetic, bitwise and
relational operator replacement, constant replacement ($k\to k\pm1$,
$c\to c+1$), unary minus deletion, \code{min}/\code{max} swap and statement
deletion. This gives \nMutAuto{} mutants (\nMutFft{} FFT, \nMutLfilter{}
\code{lfilter}, \nMutUpfirdn{} \code{upfirdn}, \nMutOla{} overlap-add).

We also use \nMutHand{} hand-written faults in six families, modelled on real
DSP bugs: a forgotten $1/N$, a twiddle sign error, single-precision
intermediates, a filter that changes after 64 samples, a gain wrong by 1\,\%, a
missing interpolation gain, a state reset between blocks. Four of them are
accuracy faults (single-precision intermediates), which first-order mutation
does not produce.

\subsection{Ground truth}
\label{sec:groundtruth}
An oracle used only for classification runs each mutant on a fixed battery of
inputs and compares it with an independent reference: all lengths up to 130
and several larger ones for the FFT; many filters, including order 0, and
initial states for \code{lfilter}; rates up to $5/4$ and $1/7$ for
\code{upfirdn}, with and without default arguments; random block splits for
overlap-add. A mutant is \emph{equivalent} if its largest relative error is at
most $10^{-12}$, an \emph{accuracy} fault if it is between $10^{-12}$ and
$10^{-6}$, and a \emph{functional} fault above $10^{-6}$, or if it raises,
returns the wrong shape or exceeds the time limit. This gives \nFunctional{}
functional faults and \nEquivalent{} equivalent mutants. No mutant falls in the
accuracy band: the largest relative error of a mutant is either below
$10^{-14}$ or above 1. First-order mutation of these subjects thus produces
large faults only, and the hand-written faults cover accuracy.

\subsection{Approaches compared}
\emph{SoP} (state of practice) is the test suite of the library whose
interface the subject implements, ported test by test with its inputs and
tolerances: NumPy's and SciPy's FFT tests, SciPy's \code{lfilter},
\code{lfilter\_zi}, \code{upfirdn} and \code{resample\_poly} tests, including
the regression test added for issue 22571. The overlap-add subject, which has no
library counterpart, gets four tests in the same style, which we wrote.
\emph{dsp-props} runs the declarations a user would write, two to four per
subject, written by the author of the relations. Three ablations isolate the components: \emph{uniform}
replaces the domain generators by uniform lengths in $[1,512]$ and values in
$[-1,1]$; \emph{fixed-tol} replaces every tolerance model by a fixed relative
tolerance of $10^{-7}$, the default of \code{numpy.testing.assert\_allclose};
\emph{generic PBT} applies both, \ie the same relations written with Hypothesis
alone.

\subsection{Metrics and execution}
For each approach and mutant we record whether the mutant is detected, the
number of generated inputs before the first failure, the length of the shrunk
counterexample and the run time. Each randomised approach runs with \nSeeds{}
seeds and we report the mean, with the range across seeds when they disagree;
SoP is deterministic and runs once. A run that exceeds \timeoutSeconds{}~s
counts as a detection, since a hanging test fails in continuous integration.
All runs use 200 inputs per relation, the library default, in the following
environment: \envDescription.

% ===========================================================================
\section{Results}
\label{sec:results}
% ===========================================================================

\subsection{RQ1: detection compared with the state of practice}
\label{sec:rq1}

\begin{table*}[t]
\caption{Mutants detected (\%, mean over seeds, with min--max when seeds disagree). For equivalent mutants the figure is a false-alarm rate and should be 0. SoP: ported test suites of NumPy and SciPy; uniform: dsp-props without its generators; fixed-tol: without its tolerance models; generic PBT: without both}
\label{tab:comparison}
\centering\footnotesize
\begin{tabular}{@{}llrrrrr@{}}
\toprule
Subject & Class (mutants) & SoP & dsp-props & uniform & fixed-tol & generic PBT \\
\midrule
FFT & functional (63) & 97 & 100 & 98 (97--100) & 100 & 99 (97--100) \\
FFT & equivalent (13) & 0 & 0 & 0 & 0 & 20 \\
\code{lfilter} & functional (51) & 100 & 86 & 86 & 86 & 86 \\
\code{lfilter} & equivalent (4) & 0 & 0 & 0 & 0 & 0 \\
\code{upfirdn} & functional (36) & 89 & 72 & 72 & 72 & 72 \\
overlap-add & functional (16) & 100 & 100 & 100 & 100 & 100 \\
\midrule
all & functional (166) & 96 & 90 & 89 (89--90) & 90 & 89 (89--90) \\
all & equivalent (17) & 0 & 0 & 0 & 0 & 15 \\
\bottomrule
\end{tabular}
\end{table*}


Table~\ref{tab:comparison} gives the share of mutants detected, averaged over
seeds. \rqOneText

\subsection{RQ2: contribution of generators and tolerances}
\label{sec:rq2}

\begin{table}[t]
\caption{Hand-written mutants by family: runs that report a failure, over mutants and 5 seeds. A correct reference must never fail}
\label{tab:hand}
\centering\scriptsize
\begin{tabular}{@{}lrrrr@{}}
\toprule
Mutant & dsp-props & uniform & fixed-tol & generic PBT \\
\midrule
FFT, reference & 0/5 & 0/5 & 0/5 & 0/5 \\
FFT, 8 mutants & 40/40 & 40/40 & 40/40 & 40/40 \\
real FFT, reference & 0/5 & 0/5 & 0/5 & 0/5 \\
real FFT, 5 mutants & 25/25 & 25/25 & 25/25 & 25/25 \\
DCT, reference & 0/5 & 0/5 & 0/5 & 0/5 \\
DCT, 3 mutants & 15/15 & 15/15 & 15/15 & 15/15 \\
IIR filter, reference & 0/5 & 0/5 & 0/5 & 0/5 \\
IIR filter, 5 mutants & 25/25 & 25/25 & 25/25 & 25/25 \\
resampler, reference & 0/5 & 0/5 & 0/5 & 0/5 \\
resampler, 6 mutants & 30/30 & 30/30 & 29/30 & 30/30 \\
streaming, reference & 0/5 & 0/5 & 0/5 & 0/5 \\
streaming, 4 mutants & 20/20 & 20/20 & 20/20 & 20/20 \\
\midrule
mutants always detected & 31 & 31 & 30 & 31 \\
first failure, median (p90) & 3 (8) & 8 (20) & 3 (13) & 8 (27) \\
counterexample length, p90 & 19 & 34 & 57 & 36 \\
\bottomrule
\end{tabular}
\end{table}


\rqTwoText

\subsection{RQ3: is the method itself correct?}
\label{sec:rq3}

A reader of any fault-finding method should ask whether its alarms are real.
We check four things.

\paragraph*{Relations} Every relation is a theorem about the declared class of
operators (Appendix~\ref{app:proofs}), so a correct implementation in exact
arithmetic never violates it; only the floating-point comparison can raise a
false alarm.

\paragraph*{False alarms} We ran every declaration on correct implementations,
NumPy's FFT and real FFT, SciPy's \code{lfilter}, \code{sosfilt},
\code{upfirdn} and \code{resample\_poly}, and the four reference subjects, with
\fpSeeds{} independent seeds: \fpRuns{} runs of a relation, \fpComparisons{}
floating-point comparisons, \fpFailed{} failures. This bounds the false-alarm
rate of these declarations on these implementations, in double precision on
one platform; it says nothing about other operators, other parameters or other
platforms, and the campaign of Section~\ref{sec:rq4} shows that a wrong
declaration does raise false alarms. SciPy's \code{lfilter} with
an initial state on empty blocks is excluded from this count because it is the
known fault of issue 22571; dsp-props flags it in \knownBugFlagged{} of
\knownBugRuns{} runs on SciPy \scipyVersion.

\paragraph*{Margin} During the same runs we recorded the ratio of the observed
error to the tolerance for every comparison. \marginExact{} of the comparisons
are exact (ratio 0); the 99th percentile is \marginPnn{} and the maximum
\marginMax: the tolerance is \marginFactor{} times larger than the worst error
observed on correct code. It is nevertheless tight enough to catch the four
single-precision faults (Table~\ref{tab:hand}).

\paragraph*{References} Calibration and the tests of the IIR model compare
against extended-precision references, a DFT computed by NumPy in
\code{long double} and a direct recursion in \code{long double}. Against
a 50-digit computation with mpmath~\cite{mpmath}, their relative error is at most
\refDftMax{}$\,\eps$ and \refLfilterMax{}$\,\eps$ respectively, where $\eps$ is
the double-precision machine epsilon: the references are three orders of
magnitude more accurate than what they are used to measure.

\paragraph*{Generators} We sampled each generator and checked its claims.
All \genPathoCount{} pathological lengths are produced, impulses have
exactly one non-zero sample, all signals are finite and of the requested
length, and amplitudes cover $10^{-100}$ to $10^{100}$; \genEmptyBlock{} of
block splits contain an empty block and \genOneBlock{} a one-sample block;
near-Nyquist frequencies lie in $[0.45,0.5]$; \genCloseBelowOne{} of the tone
pairs are less than one bin apart. The audit also found a weakness: the
``off-grid'' sinusoid lands exactly on a DFT bin in \genOnGrid{} of draws,
because Hypothesis favours simple values such as 0 and $0.5$ (\genZeroNyq{} of
draws). These are useful edge cases, but the family is less often off-grid than
its name says; a future version will draw both on-grid and strictly off-grid
frequencies.

\subsection{RQ4: faults in real libraries}
\label{sec:rq4}
We ran declarations of \campaignTargets{} operators of SciPy, NumPy, PyTorch,
librosa and torchaudio. The campaign produced the following findings.
\begin{itemize}
\item \code{scipy.signal.lfilter} returns an uninitialised final state, and
\code{sosfilt} raises, when given an empty block with an initial state. The
\code{lfilter} part is an independent rediscovery of issue 22571, fixed in the
development version after SciPy~1.18.1; the \code{sosfilt} part, which fails on
any input that is empty along the filtered axis, was new. The same fault makes
librosa's streaming pre-emphasis fail on an empty block.
\item \code{scipy.signal.tf2sos} drops pure delays: for $y[n]=x[n-2]$ it
returns an identity section, so \code{sosfilt(tf2sos(b, a), x)} differs from
\code{lfilter(b, a, x)}.
\item torchaudio's \code{deemphasis}, documented as the inverse of
\code{preemphasis}, clamps its output to $[-1,1]$.
\item PyTorch's CPU FFT reaches \torchMaxEps$\,\eps$ for $N=\torchWorstN$, against at most
\numpyMaxEps$\,\eps$ for NumPy over $N=2,\dots,1200$
(Fig.~\ref{fig:profile}); this is not a fault, but it shows why one fixed
tolerance cannot fit every library.
\end{itemize}
The campaign also produced \campaignTriaged{} alarms that triage traced to
errors in our own declarations (\eg \code{filtfilt} is not symmetric under time
reversal with its default padding); they are listed in the repository.

\begin{figure}[t]
\centering
\begin{tikzpicture}
\begin{semilogyaxis}[width=\columnwidth, height=0.55\columnwidth,
  xlabel={Length $N$}, ylabel={Worst relative error ($\eps$)},
  xmin=0, xmax=1200, ymin=0.1, ymax=400, tick label style={font=\scriptsize},
  label style={font=\footnotesize}, legend style={font=\scriptsize, draw=none, fill=none},
  legend pos=north west, legend cell align=left]
\addplot[only marks, mark=*, mark size=0.5pt, color=orange!85!black]
  table[x=N, y=torch_eps, col sep=comma]{../../benchmarks/error_profile.csv};
\addlegendentry{PyTorch (CPU)}
\addplot[only marks, mark=*, mark size=0.5pt, color=blue!70!black]
  table[x=N, y=numpy_eps, col sep=comma]{../../benchmarks/error_profile.csv};
\addlegendentry{NumPy}
\addplot[domain=2:1200, samples=100, dashed, gray] {4*ln(x)/ln(2)};
\addlegendentry{$4\log_2 N$}
\end{semilogyaxis}
\end{tikzpicture}
\caption{Worst relative FFT error, in units of $\eps$, versus length $N$, against a
\code{long double} reference (\code{benchmarks/error\_profile.py}). SciPy, not shown,
behaves like NumPy.}
\label{fig:profile}
\end{figure}

\subsection{Threats to validity}
\emph{Construct.} Mutants are a proxy for real faults~\cite{just2014mutants};
the hand-written faults and the faults of RQ4 complement them. More than half
of the detected generated mutants raise an exception or hang
(Section~\ref{sec:rq1}), which any test that runs the code would notice, so the
overall kill rates overstate what the relations themselves contribute; we
therefore also report the rates on the mutants that do not raise. The oracle
thresholds ($10^{-12}$, $10^{-6}$) are arbitrary. No mutant falls between
them, so another choice within several orders of magnitude would give the same
classification, but the per-mutant errors are in the result files so that
other thresholds can be applied.

\emph{Oracle extended after the first run.} The classification oracle was
changed once, after the approaches had been run. In the first pass it
classified \oracleFirstFunctional{} mutants as functional and
\oracleFirstEquivalent{} as equivalent; it did not call \code{upfirdn} with its
default arguments nor \code{lfilter} with an order-0 filter. The extended
oracle reclassified \nReclassified{} of the equivalent mutants as functional.
\reclDspOnly{} of them, all \reclDspOnlySubject{} mutants, were detected by
\reclDspAllVariants{} randomised variant in every seed and not by the ported
suites, so the first pass counted them as false alarms of dsp-props; the other
\reclSopOnly{}, all \reclSopOnlySubject{} mutants, were detected by the suites
only. The extension was motivated by the alarms on the first group: an oracle
refined after seeing which approach detects what is a source of bias, even
when, as here, the change moves mutants into the category where some are
missed by dsp-props. With the extended oracle, no approach other than generic
PBT detects any equivalent mutant.

\emph{Internal.} The relations, the generators, the declarations of the
comparison, the overlap-add suite and the hand-written faults were all written
by the same author, which can favour the method, in particular on the
hand-written faults. The declarations were not changed after the mutants were
first run, but the repository history does not prove that they were fixed before any mutant
was run: they were committed in the same commit as the comparison harness,
shortly before the first results, and nothing records runs made before that
commit. Results depend on random seeds, which we
vary and report.

\emph{External.} The subjects are small, and only the \code{lfilter} subject
shares the structure of its library counterpart; the SoP suites are ports of
real suites restricted to the subject's interface, and the original suites
also test other data types and array shapes. The survey counts are syntactic
lower bounds. The \campaignTargets{} targets of the campaign were chosen by
hand among the public functions of the five libraries that fit a declared
class or have a known identity (an inverse, a reference implementation, a
block decomposition), without a sampling procedure, so the campaign shows that
the method can find faults, not how often. Of these targets,
\campaignTriaged{} (\campaignTriagedShare) raised an alarm that came from an
error in our own declaration rather than from the library
(Section~\ref{sec:rq4}): declarations are themselves specifications that can be
wrong, and a user should expect to triage such alarms.

\paragraph*{Reproducing} Every number and table of this article is produced by
the following scripts, all in the \code{benchmarks/} folder of the repository:
\begin{itemize}
\item \code{survey\_test\_suites.py}: survey of Section~\ref{sec:survey};
\item \code{comparison/run\_comparison.py}: RQ1 and RQ2, generated mutants;
\item \code{comparison/run\_hand\_mutants.py}: RQ2, hand-written faults;
\item \code{validation.py}: RQ3;
\item \code{campaign.py} and \code{error\_profile.py}: RQ4.
\end{itemize}
The script \code{paper/article/make\_numbers.py} turns their result files into
the numbers and tables of the article, and \code{proofs/check.sh} checks the
Coq proofs of the appendices.

% ===========================================================================
\section{Discussion}
\label{sec:discussion}
% ===========================================================================

\paragraph*{What the relations can and cannot see} A relation detects a fault
only if the fault breaks it. Faults that change the operator into another
operator of the same class, a filter with the wrong coefficients but still LTI
and causal, a resampler that ignores the last input sample consistently,
satisfy every structural relation. The transfer-function relation closes the
first gap at the cost of one declaration; the second remains. Relations are
thus strongest on faults of structure (state, edges, indexing, normalisation)
and weakest on faults of value, where expected outputs remain useful. The
tolerances leave a gap in accuracy too: a fault that degrades a transform to a
few tens or hundreds of $\eps$, such as recursively computed twiddle factors,
passes the default constant, and \code{dp.calibrate} only refuses errors above
$10^6\eps$; a single-precision intermediate is caught, a slow loss of accuracy
is not.

\paragraph*{Parameters are part of the input} The misses of RQ1 all come from
parameter values that no declaration exercised. A practitioner should declare
an operator for several representative parameters, including degenerate ones
(order 0 and 1 filters, rate 1, default arguments), much as the SoP suites do.
Generating parameters as well as signals is possible for some classes, \eg
stable filters from random poles, and is the most direct way to extend the
method.

\paragraph*{Complementary, not competing} The method and expected-value tests
detected overlapping but different sets of faults, and together all of them;
alone, the method detected fewer. In practice they
combine: a library keeps its regression tests, which pin values and parameter
combinations, and adds relations, which need no values and explore inputs no
one thought to write down. The cost is modest: one declaration per operator and
a few seconds per run.

\paragraph*{Soundness first} A method that reports faults in other people's code
must not cry wolf. The validation of RQ3 is therefore part of the method, not an
afterthought: proofs of the relations, an error model with a measured margin,
references checked against 50-digit arithmetic, and generators audited against
their own claims. The audit found a weakness in our own generators, and the
campaign of RQ4 found \campaignTriaged{} errors in our own declarations before any fault was
reported upstream. Soundness here is measured, not proved: the false-alarm
count covers the declarations and implementations we ran, and the relations
are proved in exact arithmetic only.

% ===========================================================================
\section{Conclusion}
\label{sec:conclusion}
% ===========================================================================
Signal processing code is hard to test because the correct output is unknown,
the inputs that break it are specific, and floating-point results can only be
compared within a tolerance that nobody derives: in the libraries we surveyed,
tests compare with typed-in values, never generate their inputs, and use
tolerances chosen test by test, with \surveyScipyRelDistinct{} distinct relative
tolerances in \code{scipy.signal} alone.

We proposed a method that answers each problem from a single statement of what
an operator is: metamorphic relations proved for four classes of DSP operators
replace the expected output, signal generators target the inputs that break DSP
code, and tolerances take the form of forward error bounds, derived for linear
recursions, or come from a calibration that refuses to absorb faults. We evaluated it with a protocol that measures both its power
and its soundness. \conclusionResult{} On the correct implementations and
declarations we tested, it raised no false alarm in \fpRuns{} runs, with
references and generators that we audited.

Applied to widely used libraries, the method rediscovered a known fault of
SciPy's \code{lfilter} on empty blocks with an initial state (issue 22571,
already fixed in the development version), and found that \code{sosfilt}
raises an exception on an empty input, a new but minor edge case. It also
flagged two documented behaviours that may surprise users: \code{tf2sos} drops
pure delays, and torchaudio's \code{deemphasis} clamps its output.
\reportedUpstream

% ===========================================================================
\section*{Acknowledgments}
The author thanks Nicolas Boganski for discussions during the development of
dsp-props, and Thomas Calmant for feedback on the structure of this article
and on its experimental protocol.

\section*{AI usage disclosure}
The author conceived and directed this research, defining the problem, the
scope and the design of dsp-props, choosing the experimental protocol, checking
every result and submitting the bug reports to upstream projects. An AI coding
assistant (Claude, Anthropic) was used under the author's direction to write code and
experiment scripts and to assist with the writing. Every number in this article
is produced by a script of the repository.

\section*{Conflict of interest}
The author declares no conflict of interest.

\bibliographystyle{IEEEtran}
\bibliography{references}

\appendices

% ===========================================================================
\section{Proofs of the relations}
\label{app:proofs}
% ===========================================================================

Sequences have length $N$, indices are taken modulo $N$ for the DFT, and
$\DFT x[k]=\sum_{n=0}^{N-1}x[n]\,\omega^{kn}$ with $\omega=e^{-2\pi i/N}$.
Every proposition of this appendix is also proved in Coq, with the
Mathematical Components library, in the \code{proofs/} directory of the
repository; its README maps each proposition to its Coq theorem and lists the
modelling choices. The proofs use no axiom beyond the three classical axioms of
MathComp Analysis, which only the specialisation to $\omega=e^{-2\pi i/N}$
needs. Their scope is the algebra of the relations: they are stated in exact
arithmetic, over fields or the reals, for idealised operators. They do not
model floating-point arithmetic, do not cover the code of the library or of the
operators under test, and say nothing about the tolerances, except for the
exact-arithmetic steps of Appendix~\ref{app:iir}.

\begin{proposition}[Linearity] $\DFT$ is linear, and so is every LTI filter
from zero state and every resampler of the form ``upsample by $L$, filter,
downsample by $M$''.\end{proposition}
\begin{proof} Each is a composition of linear maps: a matrix product for
$\DFT$; zero insertion, convolution with $h$ and decimation for the resampler; for
a filter from zero state, $y$ is the convolution $h*x$ truncated to $N$ samples.
\end{proof}

\begin{proposition}[Circular shift] If $x_s[n]=x[n-s]$, then $\DFT x_s[k]=\omega^{ks}\DFT x[k]$.\end{proposition}
\begin{proof} $\sum_n x[n-s]\omega^{kn}=\sum_m x[m]\omega^{k(m+s)}$ with $m=n-s$
running over all residues.\end{proof}

\begin{proposition}[Parseval] $\|\DFT x\|_2^2=N\|x\|_2^2$.\end{proposition}
\begin{proof} $\DFT=\sqrt N\,U$ with $U$ unitary, since
$\sum_k\omega^{k(n-m)}=N\delta_{nm}$.\end{proof}
A transform declared with another normalisation states its scale, which is
checked instead of $N$.

\begin{proposition}[Hermitian symmetry] If $x$ is real, $\DFT x[N-k]=\overline{\DFT x[k]}$;
in particular $\DFT x[0]$ and, for even $N$, $\DFT x[N/2]$ are real.\end{proposition}
\begin{proof} $\omega^{(N-k)n}=\overline{\omega^{kn}}$ and $x$ is real. For
$k=0$ and $k=N/2$, the bin is its own conjugate.\end{proof}

\begin{proposition}[Inversion and convolution] $\DFT^{-1}\DFT x=x$, and
$\DFT^{-1}(\DFT x\cdot\DFT y)=x\circledast y$, the circular convolution.\end{proposition}
\begin{proof} The first follows from $\DFT^{-1}=\frac1N\overline{\DFT}$ and
Parseval's orthogonality. For the second, $\DFT(x\circledast y)[k]
=\sum_n\sum_m x[m]y[n-m]\omega^{kn}=\DFT x[k]\,\DFT y[k]$ by the shift
proposition.\end{proof}
dsp-props computes $x\circledast y$ directly in $O(N^2)$, independently of
the operator under test.

\begin{proposition}[Time invariance and causality on finite sequences] Let
$f$ be a causal LTI filter applied from zero state, $f(x)[n]=\sum_{m\le n}h[n-m]x[m]$.
Then (i) $f([0_s,x])=[0_s,f(x)]$ on $N+s$ samples, and (ii) if $x$ and $x'$
agree on $[0,k)$, then $f(x)$ and $f(x')$ agree on $[0,k)$.\end{proposition}
\begin{proof} (i) For $n<s$ the sum only involves zeros; for $n\ge s$,
substitute $m'=m-s$. (ii) $f(x)[n]$ for $n<k$ only involves $x[m]$, $m\le n<k$.
\end{proof}
For a non-causal filter (\eg a centred FIR), (i) holds after discarding the
first $s$ output samples, which is what the library checks.

\begin{proposition}[Transfer function] For a causal filter with
transfer function $H=B/A$ whose impulse response $h$ is absolutely summable,
and $k$ with $A(e^{2\pi ik/L})\neq0$, the length-$L$ DFT of
$(h[0],\dots,h[L-1])$ equals $H(e^{2\pi i k/L})$ up to the aliasing term
$\sum_{n\ge L}h[n]\omega_L^{kn}$, whose modulus is at most $\sum_{n\ge L}|h[n]|$.
\end{proposition}
\begin{proof} $H(e^{i\theta})=\sum_{n\ge0}h[n]e^{-i\theta n}$ converges
absolutely for a stable filter; split the sum at $L$.\end{proof}
The library chooses $L$ so that $r^L\approx10^{-17}$, with $r$ the largest
pole modulus, capped at $2^{16}$ samples. Since the cap leaves a large tail when
$r$ is close to 1 ($r^{L}\approx1.4\times10^{-3}$ for $r=0.9999$), the
tolerance includes a bound on $\sum_{n\ge L}|h[n]|$ rather than assuming it
small: beyond the numerator, $h$ follows the recursion $s_{n+1}=Fs_n$ with $F$
the companion matrix of $A$, and for $r<\rho<1$ the solution $P$ of
$F^{\mathsf T}PF-\rho^2P=-I$ gives
$\sum_{n\ge L}|h[n]|\le\|c\|_{P^{-1}}\|s_L\|_P/(1-\rho)$, with
$c=-a_{1:}/a_0$, whatever the multiplicity of the poles.

\begin{proposition}[Bounded gain] For an LTI filter with $\|h\|_1\le G$,
every input with $|x[n]|\le1$ gives $|y[n]|\le G$, and the input
$x[m]=\mathrm{sign}(h[N-1-m])$ attains $\sum_{n<N}|h[n]|$ at $n=N-1$.\end{proposition}
\begin{proof} $|y[n]|\le\sum_m|h[n-m]||x[m]|\le\|h\|_1$; with the given input,
every term of $y[N-1]=\sum_m h[N-1-m]x[m]$ is $|h[N-1-m]|$.\end{proof}

\begin{proposition}[Periodic time invariance] Let $L,M\ge1$ and
$f=\downarrow_M\circ(h*\cdot)\circ\uparrow_L$, the filter being applied from
zero state. Delaying the input by $Ms$
samples delays the output by $Ls$ samples.\end{proposition}
\begin{proof} $\uparrow_L$ turns a delay of $Ms$ into a delay of $LMs$; the
filter preserves it; $\downarrow_M$ turns a delay of $LMs$, a multiple of $M$,
into a delay of $Ls$.\end{proof}
The output length $\lceil NL/M\rceil$, the DC gain and the alias rejection are
declared specifications rather than theorems; their tolerances (ripple,
stop-band level) are part of the declaration.

\begin{proposition}[Block invariance] Let a streaming processor be a state
machine $(q_{t+1},y_t)=\phi(q_t,x_t)$ applied sample by sample, and let
\code{process} on a block apply $\phi$ to its samples in order, starting from the
current state and storing the final one. Then for any split
$x=x^{(1)}\cdots x^{(J)}$, including empty blocks,
$\mathrm{concat}_j\,\code{process}(x^{(j)})=\code{process}(x)$ from the same
initial state.\end{proposition}
\begin{proof} By induction on $J$: an empty block leaves the state unchanged and
outputs nothing, and processing $x^{(1)}$ then $x^{(2)}$ applies $\phi$ to the
samples of $x^{(1)}x^{(2)}$ in order.\end{proof}
The proposition holds for every split; the library checks only the splits
that its generator draws: arbitrary cut points giving up to $\max(N,16)$
blocks, including empty and one-sample ones, every block of a single sample,
or blocks of one drawn size (Section~\ref{sec:generators}).
Block-based algorithms such as overlap-add are not sample-by-sample machines,
but a correct one computes the same function, which is what the relation checks.

% ===========================================================================
\section{Error bound of a linear recursion}
\label{app:iir}
% ===========================================================================
Normalise $a_0=1$. The direct form~I recursion computes
$y[n]=\sum_{j=0}^{q}b_jx[n-j]-\sum_{j=1}^{p}a_jy[n-j]$ with $k=p+q+1$ products,
one per coefficient other than $a_0$. The bound below is derived for this
structure only; other structures, such as the transposed direct form~II, have
other error propagation paths.

\emph{Units.} Error analyses state the standard model with the unit roundoff
$u$: $\mathrm{fl}(a\mathbin{\mathrm{op}}b)=(a\mathbin{\mathrm{op}}b)(1+\theta)$
with $|\theta|\le u$, and $u=2^{-53}$ in double precision with rounding to
nearest~\cite{higham2002accuracy}. The library, like NumPy, uses the machine
epsilon $\eps=\code{np.finfo(float).eps}=2^{-52}=2u$. The derivation and the
Coq proofs below are written with $u$; we convert to $\eps$ at the end.

\emph{Derivation.} With the standard model and Higham's inner-product
bound~\cite{higham2002accuracy}, the computed $\hat y$ satisfies
$\hat y[n]=\sum_jb_jx[n-j]-\sum_ja_j\hat y[n-j]+\delta[n]$ with
$|\delta[n]|\le\gamma_k\bigl(\sum_j|b_j||x[n-j]|+\sum_j|a_j||\hat y[n-j]|\bigr)$,
$\gamma_k=ku/(1-ku)$. The Coq proof takes this bound on $\delta$ as a
hypothesis, for any $\gamma_k\ge0$; the standard model itself is not
formalised. Everything after it is proved. The error $e=\hat y-y$ obeys
$A(z)e=\delta$, \ie $e=h_A*\delta$. By Young's inequality
$\|e\|_2\le\|h_A\|_1\|\delta\|_2$, and
$\|\delta\|_2\le\gamma_k(\|b\|_1\|x\|_2+\|a_{1:}\|_1\|\hat y\|_2)$. Finally
$\|\hat y\|_2\le\|h_{B/A}\|_1\|x\|_2+\|e\|_2$, which gives, to first order and
with $\gamma_k\approx ku$,
\[
\|\hat y-y\|_2\le u\,k\,\|h_A\|_1\bigl(\|b\|_1+\|a_{1:}\|_1\|h_{B/A}\|_1\bigr)\|x\|_2 .
\]
Rigorously, if $ku\le1/2$ and $2ku\|h_A\|_1\|a_{1:}\|_1\le1/2$, the same bound
holds with $4u$ in place of $u$ for one computed output against the exact one.

\emph{Constant of the library.} In units of $\eps=2u$, the first-order bound
is \eqref{eq:iir} with $c=1/2$ and the rigorous one with $c=2$. A relation
compares computed outputs that each carry their own error, and passes as
reference norm the sum of the norms of their inputs, so the difference it
checks is bounded by twice these values: $c=1$ to first order, $c=4$
rigorously. The library uses $c=4$, \ie the rigorous bound for the difference,
under the two smallness conditions above; it ignores the few roundings of the
relation's own arithmetic, such as the scaling of the outputs in the linearity
relation. The validation of
Section~\ref{sec:rq3} confirms that the bound is never reached on correct
code, including \code{scipy.signal.lfilter}, whose transposed direct form~II
the derivation does not cover.

The factor $\|h_A\|_1$ is the worst-case peak gain (WCPG) of $1/A$, which
bounds the amplification of any bounded perturbation injected into the
recursion. The literature on the implementation of linear filters evaluates
the WCPG with guaranteed accuracy and uses it to bound rounding errors and
choose fixed-point formats~\cite{volkova2015wcpg,lopez2014formatting}; our bound is the
same argument, in $\ell_2$ norm and floating point, for the direct form~I.

\end{document}
